\chapter*{Avant-propos}
\markboth{Avant-propos}{Avant-propos}
\addcontentsline{toc}{chapter}{Avant-propos}

Ce livre étudie les inscriptions chrétiennes de Rome qui portent, ou ont reçu de leurs éditeurs, une date comprise entre 340 et 403. Il réunit \st{nA} notices. Celles-ci ne correspondent ni à autant de monuments, ni à autant de défunts : trois notices sont des doublons, deux autres sont les fragments d’une même plaque, une notice appartient au \textsc{v}\textsuperscript{e}~siècle. Une fois ces identités établies, l’enquête porte sur \st{nB} textes distincts, dont \st{nC} textes funéraires.

Le volume part d’un constat simple. Dans l’état actuel de la documentation numérique, le chercheur qui travaille sur ces épitaphes lit rarement les éditions imprimées et presque jamais les pierres. Il lit des bases de données, qui reprennent les éditions selon leurs propres conventions. Le premier objet du livre est donc cette transmission : ce qu’elle conserve, ce qu’elle perd, et ce que ses pertes font aux analyses. Le second objet est historique. Il s’agit d’examiner comment ces documents datent, nomment, qualifient les morts et inscrivent la mort dans un langage chrétien.

Les textes sont donnés d’après l’Epigraphik-Datenbank Clauss/Slaby (EDCS), dans son état de 2022. Ils ont été confrontés, notice par notice, à une seconde base savante indépendante, l’Epigraphic Database Bari (EDB), dont les fiches sont établies à partir des \emph{Inscriptiones christianae urbis Romae} et révisées par des épigraphistes nommément identifiés. Cette confrontation a établi trois doublons. Elle a aussi permis de retrouver huit formules consulaires omises par la première base, de relever les chrismes et les croix que l’EDCS ne transmet pas et de documenter, pour la première fois dans ce corpus, le cimetière d’origine, le lieu de conservation et le support de la majorité des textes. Elle n’équivaut pas à une collation des éditions imprimées : chaque notice indique ce qui a été contrôlé, et comment.

Cette double exigence commande la forme du livre. Chaque proposition générale renvoie à des notices identifiées ; chaque compte annonce son dénominateur et sa règle ; chaque restitution moderne reste distinguée des lettres lues. Lorsque deux lectures publiées s’opposent, les deux sont présentées et l’année de classement n’est pas modifiée sans preuve.

\vfill
\noindent\emph{Brive-la-Gaillarde, le 14 octobre 2026}\hfill A. A. M.

\chapter*{Résumé}
\markboth{Résumé · Abstract}{Résumé · Abstract}
\addcontentsline{toc}{chapter}{Résumé · Abstract · Mots-clés}

Le volume édite et analyse \st{nA} notices d’inscriptions chrétiennes de Rome datées de 340 à 403. La confrontation systématique des textes de l’Epigraphik-Datenbank Clauss/Slaby (EDCS) avec l’Epigraphic Database Bari (EDB) porte sur \st{edbApp} notices sur \st{nA}. Elle fusionne trois doublons (106 dans 120, 282 dans 288, 329 dans 336), réunit deux fragments d’une même plaque (393 et 394) et écarte un texte du \textsc{v}\textsuperscript{e}~siècle (\nnum{98}, daté de 461 ou de 482). Le corpus analysé compte ainsi \st{nB} textes distincts, dont \st{nC} textes funéraires. La seconde base transmet huit formules consulaires absentes de la première, une partie grecque omise et, sur \st{edbAppC} textes funéraires appariés, \st{nSignesC} chrismes ou croix (\st{pSignesC}~\%) que l’EDCS ne signale pas. L’année de \st{datLu} textes funéraires repose sur un nom consulaire lu, celle de \st{datPartiel} sur une lecture partielle, celle de \st{datRest} sur un consulat restitué ; dix datations restent contestées. La formule \emph{in pace} est attestée dans \st{ip_tot} textes funéraires (\st{ip_pct}~\%), dont \st{ip_ecrit} en toutes lettres et \st{ip_restitue} seulement par restitution. Les désignations écrites de Dieu, du Christ, de l’Esprit ou de la Trinité concernent \st{nProfilLarge} textes (\st{pProfilLarge}~\%), et \st{nProfilStrict} seulement (\st{pProfilStrict}~\%) si l’on écarte les lectures dépendant d’une restitution ou d’un monogramme. Aucune mention de la résurrection n’est assurée : la seule attestation alléguée est une restitution contestée. Les comptes décrivent une documentation transmise, non la population chrétienne de Rome.

\section*{Abstract}
This volume edits and analyses \st{nA} entries for Christian inscriptions from Rome dated between AD 340 and 403. Texts transmitted by the Epigraphik-Datenbank Clauss/Slaby (EDCS) were systematically compared with the Epigraphic Database Bari (EDB) for \st{edbApp} of the \st{nA} entries. The comparison identifies three duplicates, joins two fragments of one slab and removes one fifth-century text (no.~98, AD 461 or 482). The analysed corpus therefore comprises \st{nB} distinct texts, \st{nC} of them funerary. The second database supplies eight consular formulae missing from the first, an omitted Greek section and, among \st{edbAppC} matched funerary texts, \st{nSignesC} christograms or crosses (\st{pSignesC}\%) not recorded by EDCS. The year of \st{datLu} funerary texts rests on a consular name that is actually read; ten datings remain disputed. \emph{In pace} occurs in \st{ip_tot} funerary texts (\st{ip_pct}\%), \st{ip_ecrit} of them written in full. Written designations of God, Christ, the Spirit or the Trinity occur in \st{nProfilLarge} texts (\st{pProfilLarge}\%), and in only \st{nProfilStrict} (\st{pProfilStrict}\%) once readings depending on restorations or monograms are excluded. No mention of the resurrection is secure. The figures describe a transmitted documentation, not Rome’s Christian population.

\section*{Mots-clés · Keywords}
Rome ; Antiquité tardive ; épigraphie chrétienne ; épitaphes ; datation consulaire ; chronologie ; formulaire funéraire ; humanités numériques ; EDCS ; EDB ; critique documentaire.

\noindent Rome; Late Antiquity; Christian epigraphy; epitaphs; consular dating; chronology; funerary formulary; digital epigraphy; EDCS; EDB; source criticism.

\chapter*{Sigles et abréviations}
\markboth{Sigles et abréviations}{Sigles et abréviations}
\addcontentsline{toc}{chapter}{Sigles et abréviations}
\begin{tabularx}{\linewidth}{@{}l X@{}}
AE & \emph{L’Année épigraphique} \\
CIL & \emph{Corpus Inscriptionum Latinarum} \\
CLE & \emph{Carmina Latina Epigraphica} (F. Bücheler, E. Lommatzsch) \\
CLRE & R. S. Bagnall, A. Cameron, S. R. Schwartz, K. A. Worp, \emph{Consuls of the Later Roman Empire} (1987) \\
EDB & Epigraphic Database Bari (université de Bari) \\
EDCS & Epigraphik-Datenbank Clauss/Slaby \\
EDH & Epigraphische Datenbank Heidelberg \\
EDR & Epigraphic Database Roma \\
ICUR, ICVR & \emph{Inscriptiones christianae urbis Romae} ; « ICUR n.s. » ou « ICUR » suivi d’un chiffre romain : \emph{nova series} (1922–1992) ; « ICUR I » de De Rossi (1857–1861) précisé lorsque nécessaire \\
ILCV & E. Diehl, \emph{Inscriptiones Latinae Christianae Veteres} \\
InscrIt & \emph{Inscriptiones Italiae} \\
PCAS & Pontificia Commissione di Archeologia Sacra (photographies signalées par l’EDB) \\
PCBE & \emph{Prosopographie chrétienne du Bas-Empire} \\
PLRE & \emph{The Prosopography of the Later Roman Empire} \\
SDAM & Social Dynamics and Complexity in the Ancient Mediterranean (université d’Aarhus) \\
SEG & \emph{Supplementum epigraphicum Graecum} \\
\midrule
A, B, C & ensembles analytiques définis au chapitre~\ref{ch:corpus} \\
p.c. & \emph{post consulatum} \\
v.c. & \emph{vir clarissimus}, \emph{viri clarissimi} \\
\end{tabularx}
